\chapter{Real Sheaves and Purity}
\label{chap:real-sheaves}

The theorem of this chapter is the engine that prevents the Fourier
argument from being circular.  It extracts punctual purity from a
reality condition by using positivity of even tensor powers.  Its proof
uses the trace formula on a curve, the rank-one theorem of
\cref{found:rank-one-purity}, elementary monodromy, and Poincare duality
in top degree.  It does \emph{not} use the direct-image weight theorem
which it will later help prove.

Fix an isomorphism \(\iota:\Qbarl\to\mathbf C\) throughout this chapter.

\section{Reality and positivity}

\begin{definition}
\label{def:iota-real}
A constructible Weil sheaf \(\cG_0\) on \(X_0\) is
\emph{\(\iota\)-real} if, for every closed point \(x_0\),
\[
  \iota\det(1-TF_{x_0}\mid\cG_{\bar x})
  \in\mathbf R[T].
\]
Equivalently, the multiset of complex numbers obtained from the local
Frobenius eigenvalues is stable under complex conjugation.
\end{definition}

Reality is preserved by direct sums, duals, tensor powers, exterior
powers, subquotients whose complementary factor is also real, and
semisimplification.  It is not in general inherited by an individual
irreducible constituent.  This last observation is why the proof below
needs exterior powers rather than simply applying the irreducible case
constituent by constituent.

\begin{lemma}[Positive local factors]
\label{lem:positive-local-factors}
If \(\cG_0\) is \(\iota\)-real and \(m\geq 1\), then for every closed
point \(x_0\) the series
\[
  \iota\det\left(1-TF_{x_0}\mid
  \cG_{\bar x}^{\otimes 2m}\right)^{-1}
\]
has nonnegative real coefficients.
\end{lemma}

\begin{proof}
For an endomorphism \(A\) of a finite-dimensional vector space,
\[
  \log\det(1-TA)^{-1}
  =\sum_{n\geq1}\Tr(A^n)\frac{T^n}{n}.
\]
Since the characteristic polynomial of \(F_{x_0}\) becomes real under
\(\iota\), all power traces
\(\iota\Tr(F_{x_0}^n\mid\cG_{\bar x})\) are real.  On the tensor power,
\[
  \Tr(F_{x_0}^n\mid\cG_{\bar x}^{\otimes 2m})
  =\Tr(F_{x_0}^n\mid\cG_{\bar x})^{2m},
\]
which becomes a nonnegative real number.  The exponential of a formal
series with zero constant term and nonnegative coefficients again has
nonnegative coefficients.
\end{proof}

We shall use one elementary fact about positive products.

\begin{lemma}[A local factor cannot converge less far]
\label{lem:positive-product-radius}
Let \(f_j(T)\in1+T\mathbf R_{\geq0}[[T]]\) be a locally finite family in
the \(T\)-adic topology, and put \(f=\prod_j f_j\).  The radius of
convergence of each \(f_j\) is at least the radius of convergence of
\(f\).
\end{lemma}

\begin{proof}
Write \(f_j=\sum_na_{j,n}T^n\) and \(f=\sum_na_nT^n\).  On multiplying
series with nonnegative coefficients, the contribution obtained by
taking the degree \(n\) term from \(f_j\) and constant terms from all
other factors shows \(0\leq a_{j,n}\leq a_n\).  The
Cauchy--Hadamard formula gives
\[
  \limsup_n a_{j,n}^{1/n}
  \leq\limsup_n a_n^{1/n},
\]
which is the asserted comparison of radii.
\end{proof}

\section{Top cohomology of an affine curve}

Let \(C_0\) be a smooth geometrically connected affine curve over \(k_0\),
and let \(\cG_0\) be a lisse Weil sheaf corresponding to a representation
\(V\) of \(W(C_0,\bar c)\).

\begin{lemma}
\label{lem:curve-hc0-hc2}
There are Frobenius-equivariant identifications
\[
  H_c^0(C,\cG)=0,\qquad
  H_c^2(C,\cG)\simeq
  V_{\pi_1(C,\bar c)}(-1),
\]
where the subscript denotes geometric coinvariants.  In particular, if
\(\lambda\) is an eigenvalue of geometric Frobenius on \(H_c^2(C,\cG)\),
then
\[
  2\log_q|\iota(\lambda/q)|
\]
is a determinant weight of a geometrically constant quotient of
\(\cG_0\), and hence is among the determinant weights of \(\cG_0\).
\end{lemma}

\begin{proof}
A lisse section with proper support vanishes on a nonempty open set and
therefore vanishes everywhere on the connected curve, proving the first
claim.  Poincare duality gives
\[
 H_c^2(C,\cG)
 \simeq H^0(C,\cG^\vee)^\vee(-1).
\]
The space \(H^0(C,\cG^\vee)\) is
\((V^\vee)^{\pi_1(C,\bar c)}\); its dual is
\(V_{\pi_1(C,\bar c)}\).  The twist \((-1)\) multiplies geometric
Frobenius by \(q\).  The coinvariant representation is geometrically
constant and is a quotient of \(V\).  Its irreducible Weil constituents
are rank one after extension of scalars because the geometric group acts
trivially.  Their weights are determinant weights of constituents of
\(\cG_0\), by \cref{found:rank-one-purity,def:determinant-weight}.
\end{proof}

\begin{corollary}
\label{cor:hc2-upper-det-weight}
If \(\beta\) is the largest determinant weight of a geometrically
semisimple \(\cG_0\), every eigenvalue \(\lambda\) on
\(H_c^2(C,\cG)\) satisfies
\[
  |\iota(\lambda)|\leq q^{(\beta+2)/2}.
\]
For \(\cG_0^{\otimes r}\), the same assertion holds with \(\beta\)
replaced by \(r\beta\).
\end{corollary}

\begin{proof}
The first statement is immediate from
\cref{lem:curve-hc0-hc2}; the second uses
\cref{lem:det-weight-arithmetic}.
\end{proof}

\section{The affine-curve real-sheaf theorem}

\begin{theorem}[Real sheaves on affine curves]
\label{thm:real-sheaf-affine-curve}
Let \(C_0\) be a smooth geometrically connected affine curve, and let
\(\cG_0\) be a lisse \(\iota\)-real Weil sheaf.  Every irreducible
constituent of its semisimplification is punctually
\(\iota\)-pure, of weight equal to its determinant weight.  Consequently
\(\cG_0\) is \(\iota\)-mixed.
\end{theorem}

\begin{proof}
Semisimplification leaves all local characteristic polynomials unchanged,
so we may assume that \(\cG_0\) is semisimple.

\medskip
\noindent\emph{Step 1: a universal upper bound.}
Let \(\beta\) be the largest determinant weight of \(\cG_0\), and fix
\(m\geq1\).  The trace formula for the even tensor power, together with
\cref{lem:curve-hc0-hc2}, gives an identity of rational functions and of
Taylor series near \(0\):
\[
\begin{split}
 L(C_0,\cG_0^{\otimes 2m},T)
 &=
 \prod_{x_0\in|C_0|}
 \det\left(1-T^{\deg(x_0)}F_{x_0}\mid
 \cG_{\bar x}^{\otimes 2m}\right)^{-1}\\
 &=
 \frac{
 \det(1-TF\mid H_c^1(C,\cG^{\otimes 2m}))
 }{
 \det(1-TF\mid H_c^2(C,\cG^{\otimes 2m}))
 }.
\end{split}
\tag{9.1}
\label{eq:real-sheaf-L}
\]
By \cref{cor:hc2-upper-det-weight}, every zero of the denominator on the
right has modulus at least
\[
  q^{-(2m\beta+2)/2}=q^{-m\beta-1}.
\]
The numerator can cancel a pole but cannot create a new one.  Hence the
Taylor series of the rational function in \eqref{eq:real-sheaf-L} has
radius of convergence at least \(q^{-m\beta-1}\).

All local factors on the first line have nonnegative coefficients after
applying \(\iota\), by \cref{lem:positive-local-factors}.  The product is
locally finite because only finitely many closed points have bounded
degree.  By \cref{lem:positive-product-radius}, each local factor has
radius at least \(q^{-m\beta-1}\).

Let \(\alpha\) be an eigenvalue of \(F_{x_0}\) on
\(\cG_{\bar x}\), with \(e=\deg(x_0)\).  Then \(\alpha^{2m}\) is an
eigenvalue on the even tensor power.  The corresponding local reciprocal
factor has a singularity at a complex number \(T\) satisfying
\[
  |T|^e=|\iota(\alpha)|^{-2m}.
\]
The radius estimate therefore gives
\[
  |\iota(\alpha)|^{2m/e}
  \leq q^{m\beta+1},
\]
or
\[
  |\iota(\alpha)|
  \leq q^{e(\beta/2+1/(2m))}.
\]
Letting \(m\) tend to infinity yields the universal bound
\[
  |\iota(\alpha)|\leq q^{e\beta/2}.
\tag{9.2}
\label{eq:universal-real-upper}
\]
Notice that no estimate for \(H_c^1\) was used: it occurs only in the
numerator of \eqref{eq:real-sheaf-L}.  This is the essential
noncircularity of the argument.

\medskip
\noindent\emph{Step 2: the irreducible case.}
Suppose first that \(\cG_0\) is irreducible, so its sole determinant
weight is \(\beta\).  Its dual is again \(\iota\)-real and has determinant
weight \(-\beta\).  Applying \eqref{eq:universal-real-upper} to the
eigenvalue \(\alpha^{-1}\) of the dual gives
\[
  |\iota(\alpha)|^{-1}\leq q^{-e\beta/2}.
\]
Together with \eqref{eq:universal-real-upper}, this proves
\[
  |\iota(\alpha)|=q^{e\beta/2}
\]
at every closed point.  Thus \(\cG_0\) is pure of its determinant weight.

\medskip
\noindent\emph{Step 3: separation by exterior powers.}
An individual constituent of a real sheaf need not be real, so Step 2
cannot simply be applied to each constituent.  Group the irreducible
constituents by strictly decreasing determinant weights
\[
  \gamma_1>\gamma_2>\cdots>\gamma_s,
\]
and let \(r_j\) be the total rank of the constituents of weight
\(\gamma_j\).  Put \(N_j=r_1+\cdots+r_{j-1}\).

The sheaf \(\bigwedge^{N_j+1}\cG_0\) is \(\iota\)-real: its local
characteristic polynomial is a universal polynomial with integer
coefficients in that of \(\cG_0\).  The largest determinant weight of
this exterior power is
\[
  r_1\gamma_1+\cdots+r_{j-1}\gamma_{j-1}+\gamma_j.
\tag{9.3}
\label{eq:exterior-largest}
\]
Indeed, the largest central weight is obtained by taking every available
direction in the higher-weight blocks and one direction in the
\(\gamma_j\)-block.

Fix a closed point of degree \(e\) and an eigenvalue \(\alpha\) belonging
to a constituent of determinant weight \(\gamma_j\).  Multiply it by all
the Frobenius eigenvalues, with multiplicity, in the blocks of weights
\(\gamma_1,\ldots,\gamma_{j-1}\).  The resulting number \(\delta\) is an
eigenvalue on \(\bigwedge^{N_j+1}\cG_0\).  The absolute value of the
product contributed by the higher blocks is known even before their
individual purity: the determinant definition gives
\[
 \left|\iota\prod_{\text{higher blocks}}\alpha_h\right|
 =
 q^{\frac e2(r_1\gamma_1+\cdots+r_{j-1}\gamma_{j-1})}.
\]
Apply the universal upper bound to the exterior power and use
\eqref{eq:exterior-largest}.  Cancellation of the known determinant
factor yields
\[
  |\iota(\alpha)|\leq q^{e\gamma_j/2}.
\tag{9.4}
\]
Now repeat the exterior-power argument for \(\cG_0^\vee\).  Its weight
blocks occur in reverse order and the eigenvalue corresponding to
\(\alpha\) is \(\alpha^{-1}\), of determinant weight \(-\gamma_j\).
The resulting upper bound for \(\alpha^{-1}\) is the reverse inequality
to \eqref{eq:universal-real-upper}.  Hence
\[
  |\iota(\alpha)|=q^{e\gamma_j/2}.
\]
Every constituent is pure of its determinant weight.  A Jordan--Holder
filtration of the original sheaf is therefore a weight filtration, so
the possibly nonsemisimple \(\cG_0\) is mixed.
\end{proof}

\begin{remark}
The tensor-power limit in this proof is a local positivity argument, not
the Kunneth tensor-power argument from Deligne's Weil I proof.  No
Lefschetz pencil, Chebotarev theorem, or coarse middle-cohomology estimate
appears.  The two arguments should not be conflated.
\end{remark}

\section{From curves to arbitrary dimension}

We require a Bertini statement that preserves irreducibility of a given
local system, not merely smoothness of the curve.

\begin{foundation}
\label{found:irreducibility-bertini}
\textbf{Bertini irreducibility capsule.}
Let \(X\) be smooth and geometrically irreducible over an algebraically
closed field, and let \(\cF\) be an irreducible lisse
\(\Qbarl\)-sheaf.  After choosing a sufficiently ample embedding, the
linear sections \(C=X\cap L\) which are smooth geometrically irreducible
curves and for which \(\cF|_C\) remains irreducible form a nonempty open
subset of the relevant Grassmannian.

The geometric input is proved by replacing the representation by a finite
collection of finite quotients that detect invariant subspaces, applying
the fundamental-group Bertini theorem to obtain surjectivity onto those
quotients, and using constructibility of the exceptional locus.  A
compactness argument in the profinite fundamental group passes from the
finite detectors to the \(\ell\)-adic representation.  Over a finite
field, a point of the nonempty Grassmannian open exists after a finite
extension; taking all Galois conjugates descends a dense open on which the
slicing conclusion holds.  This is independent of weights; see
\cite[Appendix~B]{KiehlWeissauer} and \cite{StanfordL22}.
\end{foundation}

\begin{theorem}[Real sheaves are mixed]
\label{thm:real-sheaf-mixed}
Let \(X_0\) be a separated finite-type \(k_0\)-scheme and let
\(\cG_0\) be a constructible \(\iota\)-real Weil sheaf.
\begin{enumerate}
\item The sheaf \(\cG_0\) is \(\iota\)-mixed.
\item If \(X_0\) is normal and geometrically irreducible and \(\cG_0\)
is lisse, every irreducible constituent is punctually
\(\iota\)-pure of its determinant weight.
\end{enumerate}
\end{theorem}

\begin{proof}
We use Noetherian induction on the support.  Nilpotents do not affect the
\'etale topos, so assume \(X_0\) reduced.  Choose a dense open immersion
\(j_0:U_0\hookrightarrow X_0\) such that \(U_0\) is smooth and every
cohomology sheaf under consideration is lisse.  With \(i_0\) the reduced
closed complement, there is an exact sequence
\[
  0\longrightarrow j_{0!}j_0^*\cG_0
  \longrightarrow\cG_0
  \longrightarrow i_{0*}i_0^*\cG_0
  \longrightarrow0.
\tag{9.5}
\label{eq:real-sheaf-stratification}
\]
The closed term is mixed by induction.  It remains to prove purity of the
irreducible constituents on a dense open of each irreducible component.
After a finite extension of \(k_0\) we may suppose the component
geometrically irreducible and the constituents absolutely irreducible;
mixedness descends because the original sheaf is a direct summand of
finite pushforward of its pullback.

Let \(\cF_0\) be one such constituent.  If the dimension is one,
\cref{thm:real-sheaf-affine-curve} applies after deleting finitely many
points.  In higher dimension, use
\cref{found:irreducibility-bertini}.  For every geometric point in a
suitable dense open there is a good curve through that point, defined
after a finite extension, on which \(\cF\) remains irreducible.  The
restriction of the entire \(\cG_0\) is still real, so
\cref{thm:real-sheaf-affine-curve} makes \(\cF|_C\) pure.  By
\cref{cor:dense-pullback-det-weights}, its weight is the determinant
weight of \(\cF_0\), independent of the chosen curve.  Thus all closed
points in the dense open have the required punctual weight.

The Jordan--Holder constituents of \(j_0^*\cG_0\) are therefore pure.
They give a mixed filtration on \(U_0\); extension by zero preserves the
same local Frobenius polynomials on \(U_0\) and has zero stalks outside.
Combining that filtration with
\eqref{eq:real-sheaf-stratification} and the induction filtration on the
complement proves (1).  The same argument without
stratification proves (2).
\end{proof}

\section{Real envelopes}

The Fourier transform of a pure sheaf need not visibly be pure.  We shall
first place it inside a real sheaf and invoke
\cref{thm:real-sheaf-mixed}.

\begin{lemma}[Real envelope of a pure sheaf]
\label{lem:real-envelope}
Let \(\cF_0\) be punctually \(\iota\)-pure of weight \(w\).  There is a
rank-one geometrically constant Weil sheaf \(\cL_b\), with degree-one
Frobenius scalar \(b\), such that
\[
  \cF_0\oplus(\cF_0^\vee\otimes\cL_b)
\]
is \(\iota\)-real and contains \(\cF_0\) as a direct summand.
\end{lemma}

\begin{proof}
Choose \(b\in\Qbarl^\times\) with
\(\iota(b)=q^w\); the existence uses only that \(\iota\) is a field
isomorphism.  Let \(\cL_b\) be the Weil line on which a degree \(n\)
Frobenius acts by \(b^n\).  If \(\alpha\) is a local eigenvalue at a
degree \(n\) point, purity gives
\[
 \overline{\iota(\alpha)}
 =\frac{q^{nw}}{\iota(\alpha)}
 =\iota(b^n/\alpha).
\]
Thus the eigenvalues on \(\cF_0^\vee\otimes\cL_b\) are precisely the
complex conjugates of those on \(\cF_0\), with multiplicity.  The
characteristic polynomial of their direct sum is real.
\end{proof}

\begin{warning}
The line \(\cL_b\) is generally only a Weil sheaf; \(b\) need not define a
continuous representation of \(\widehat{\mathbf Z}\).  This is exactly
why \cref{warn:weil-versus-etale} was necessary.
\end{warning}

\begin{formalization}
\label{form:real-sheaf-interface}
\textbf{Interface exported by the real-sheaf chapter.}

\begin{enumerate}
\item The predicate \(\operatorname{Real}_\iota(\cG)\), characterized
equivalently by real power traces at every closed point.
\item Positivity of the inverse local determinant for even tensor powers.
\item The positive-product radius lemma as a statement about
nonnegative formal power series.
\item The curve calculations \(H_c^0=0\) and
\(H_c^2\simeq V_{\pi_1}(-1)\), Frobenius-equivariantly.
\item The universal determinant-weight upper bound and its limiting
argument.
\item Exterior-power separation of determinant-weight blocks.
\item Bertini slicing as a separately named geometric capability.
\item The terminal theorem
\(\operatorname{Real}_\iota(\cG)\Rightarrow
\operatorname{Mixed}_\iota(\cG)\).
\item A real-envelope constructor for a pure Weil sheaf.
\end{enumerate}

The dependency audit for the terminal theorem contains trace formula,
top-degree duality, rank-one purity, algebraic monodromy, and Bertini.  It
does not contain the curve estimate of \cref{chap:curve-estimate}, Fourier
transform, or the direct-image theorem of
\cref{chap:global-weights}.
\end{formalization}



